% BOOK-FILE-ID: f-dcd823eec8b6
% section: 数学的帰納法

\section{数学的帰納法}\label{節：数学的帰納法}
数項からの予想を,すべての項の証明へ進める方法である.
自然数$n$の主張$P(n)$について,次の二つを示す.
\begin{enumerate}[label=(\roman*),parsep=3pt]
  \item $P(1)$が成り立つ.
  \item 任意の$n\ge1$について,$P(n)$が成り立つなら$P(n+1)$も成り立つ.
\end{enumerate}
(i)から出発し,(ii)を繰り返すと,$P(1)\Rightarrow P(2)\Rightarrow P(3)\Rightarrow\cdots$となる.
ドミノ倒しなら,(i)は最初の一枚を倒すこと,(ii)はどの一枚からも次へ伝わる並べ方に当たる.

\sectionref*{節：数列と漸化式}で確かめた$a_n=2n-1$では,初項の一致が(i),
$a_k=2k-1$から漸化式で$a_{k+1}=2(k+1)-1$を得た計算が(ii)である.
すべての$n$で正しいと仮定するのではなく,\textbf{ある番号で正しければ次でも正しい}ことを示す.

\sectionref*{節：複素平面と極形式}の\bookterm{de-moivre}{ド・モアブルの定理}や,
\sectionref*{節：特性方程式の解が虚数の場合}の\bookterm{general}{一般項}もこの方法で確かめる.

\begin{bookpoint}
  \textbf{POINT}\par
  出発点での成立と,ある番号での成立から次の番号での成立が従うことを示す.
  漸化式では,必要な初期値を確かめた上で,与えられた規則を使って次の項へ進める.
\end{bookpoint}
